\subsection{无穷次可微函数的 Fourier 变换}

回忆 (Chap. III, §3, no. 1, Def. 2)，U(G) 表示支集包含于 \(\{e\}\) 的 G 上分布所成的代数。g 到 U(G) 的典则嵌入延拓为从李代数 g 的包络代数到 U(G) 的同构 (loc. cit., no. 7, Prop. 25)；以下通过此同构识别这两个代数。若 f 是 G 上的无穷次可微复函数，且 \(t \in \mathrm{U}(\mathrm{G})\)，记 \(L_{t}f\) 和 \(R_{t}f\) 为 G 上由下式定义的函数

\[
\mathrm{L} _ {t} f (g) = \langle \varepsilon_ {g} * t, f \rangle , \quad \mathrm{R} _ {t} f (g) = \langle t * \varepsilon_ {g}, f \rangle
\]

(cf.~loc. cit., no. 6). 对于所有 \(g \in G\)，

\[
\mathrm{L} _ {t} \circ \gamma (g) = \gamma (g) \circ \mathrm{L} _ {t}, \quad \mathrm{R} _ {t} \circ \delta (g) = \delta (g) \circ \mathrm{R} _ {t}.
\]

设 \(u \in \hat{G}\)；记 \(E_u\) 为 \(u\) 的表示空间。李群态射 \(u: G \to \mathbf{GL}(E_u)\) 经微分给出实李代数的同态 \(\mathfrak{g} \to \text{End}(E_u)\)，从而给出代数同态，仍记作 \(u\)，即从 \(U(G)\) 到 \(\text{End}(E_u)\) 的同态。若 \(t \in U(G)\) 且 \(f\) 是 \(G\) 上的无穷次可微函数，则

\[
u (\mathrm{L} _ {t} f) = u (f) u (t ^ {\vee}), \quad u (\mathrm{R} _ {t} f) = u (t ^ {\vee}) u (f),\tag{14}
\]

其中 \(t^{\vee}\) 表示 t 在 U(G) 的主反自同构下的像 (Chap. I, §2, no. 4)；事实上，只需对 \(t \in g\) 验证此式，而此时它由公式 (12) 和 (13) 微分得到 (cf.~Chap. III, §3, no. 7, Prop. 27)。

对于所有 \(u \in \hat{\mathbf{G}}\)，记 \(\lambda(u)\) 为 \(u\) 的最高权 (§7, no. 2, Th. 1)，于是 \(u \mapsto \lambda(u)\) 是从 \(\hat{\mathbf{G}}\) 到由 \(\mathrm{X}_{++}\) 的支配元组成的集合 \(\mathrm{X}(\mathrm{T})\) 的双射。

设 \(\Gamma \in \mathrm{U}(\mathrm{G})\) 是 \(\mathrm{G}\) 的 Casimir 元 (§7, no. 6)；对于所有 \(u \in \hat{\mathrm{G}}\)，\(u(\Gamma)\) 是 \(\mathrm{E}_u\) 上的数乘变换，其倍数记为 \(\tilde{\Gamma}(u)\)，因此得到映射 \(u \mapsto \tilde{\Gamma}(u)\)，从 \(\hat{\mathrm{G}}\) 到 \(\mathbf{C}\)。

若 \(\varphi\) 和 \(\psi\) 是 \(\hat{G}\) 上取正实值的两个函数，则用 `` \(\varphi \preccurlyeq \psi\) '' 或 `` \(\varphi(u) \preccurlyeq \psi(u)\) '' 表示关系 ``存在 M \textgreater{} 0，使得 \(\varphi(u) \leq M\psi(u)\) 对所有 \(u \in \hat{G}\) 都成立''；这是 \(\hat{G}\) 上取正实值的函数集上的预序关系。

命题 2. 设 \(m \mapsto \|m\|\) 是实向量空间 \(\mathbf{R} \otimes \mathrm{X(T)}\) 上的一个范数，且 \(\Gamma\) 是 G 的 Casimir 元。设 \(\varphi\) 是 \(\hat{\mathrm{G}}\) 上取正实值的函数。

\begin{enumerate}
\def\labelenumi{\alph{enumi})}
\tightlist
\item
  下列条件等价：
\end{enumerate}

\begin{enumerate}
\def\labelenumi{(\roman{enumi})}
\item
  存在整数 \(n > 0\)，使得 \(\varphi(u) \preccurlyeq (\|\lambda(u)\| + 1)^n\)（相应地，对于每个整数 \(n > 0\)，有 \(\varphi(u) \preccurlyeq (\|\lambda(u)\| + 1)^{-n}\)）。
\item
  存在整数 \(n > 0\)，使得 \(\varphi(u) \preccurlyeq (\tilde{\Gamma}(u) + 1)^n\)（相应地，对于每个整数 \(n > 0\)，有 \(\varphi(u) \preccurlyeq (\tilde{\Gamma}(u) + 1)^{-n}\)）。
\end{enumerate}

\begin{enumerate}
\def\labelenumi{\alph{enumi})}
\setcounter{enumi}{1}
\tightlist
\item
  若 \(\mathrm{G}\) 是半单的，则条件 (i) 和 (ii) 还等价于：
\end{enumerate}

\begin{enumerate}
\def\labelenumi{(\roman{enumi})}
\setcounter{enumi}{2}
\tightlist
\item
  存在整数 \(n > 0\)，使得 \(\varphi(u) \preccurlyeq d(u)^n\)（相应地，对于每个整数 \(n > 0\)，有 \(\varphi(u) \preccurlyeq d(u)^{-n}\)）。
\end{enumerate}

首先注意，条件 (i) 显然与范数的选择无关。因此可以使用由二次型 \(Q_{\varGamma}\)（它与 \(\varGamma\) 相关）定义的范数 (§7, no. 6, Prop. 4)。于是

\[
0 \leq \tilde {\Gamma} (u) = \| \lambda (u) + \rho \| ^ {2} - \| \rho \| ^ {2},
\]

所以 \(\tilde{\Gamma}(u) + 1 \preccurlyeq (\|\lambda(u)\| + 1)^2 \preccurlyeq \tilde{\Gamma}(u) + 1\)，故得 \(a\)。

此外，若 G 是半单的，

\[
\left\| \lambda (u) + \rho \right\| \preccurlyeq d (u) \preccurlyeq \left\| \lambda (u) + \rho \right\| ^ {\mathrm{N}}, \quad \text { where } \mathrm{N} = 1 / 2 (\dim \mathrm{G} - \dim \mathrm{T})
\]

(§7, no. 5, Cor. 1 of Th. 3)，所以 \(\| \lambda(u) \| + 1 \preccurlyeq d(u) \preccurlyeq (\| \lambda(u) \| + 1)^{\mathrm{N}}\)，故得 \(b\)）。

由命题 2 可知，条件 (i) 与极大环面、房和范数的选择无关，而条件 (ii) 与 Casimir 元的选择无关。函数 \(\varphi\) 若满足条件 (i) 和 (ii)，称为缓增函数（相应地，速降函数）。两个缓增函数的乘积是缓增的；缓增函数与速降函数的乘积是速降的。若 \(\varphi\) 是速降函数，则族 \((\varphi(u))_{u\in\hat{G}}\) 可求和。

例。函数 \(u \mapsto d(u)\) 是缓增的 (§7, no. 5, Cor. 1 of Th. 3)；对于任意范数 \(\| \cdot \|\)，其定义在 \(\mathbf{R} \otimes \mathrm{X(T)}\) 上，函数 \(u \mapsto \| \lambda(u) \|\) 是缓增的。对于任意 Casimir 元 \(\Gamma\)，函数 \(u \mapsto \tilde{\Gamma}(u)\) 是缓增的；更一般地：

命题 3. 对于所有 \(t \in \mathrm{U}(\mathrm{G})\)，函数 \(u \mapsto \| u(t) \|_{\infty}\) 和 \(u \mapsto \| u(t) \|_2\) 在 \(\hat{\mathrm{G}}\) 上都是缓增的。

由于两个缓增函数的乘积仍为缓增的，只需在 \(t \in g\) 时证明此断言；在这种情形下，断言由 §7, no. 6 中的 Remark 以及不等式

\[
\| u (t) \| _ {2} \leq d (u) \| u (t) \| _ {\infty}.
\]

定理 1. a) 设 \(f\) 是 \(\mathbf{G}\) 上的无穷次可微复函数。则族 \((f_u)_{u\in \hat{\mathbf{G}}}\)，其中 \(f_{u}(g) = \langle u(g)|u(f)\rangle\)，在 \(\mathbf{G}\) 上一致可求和，并且对于所有 \(h\in \mathbf{G}\)，

\[
f (h) = \sum_ {u \in \hat {\mathrm{G}}} \langle u (h) | u (f) \rangle = \sum_ {u \in \hat {\mathrm{G}}} d (u) \int_ {\mathrm{G}} f (g) \mathrm{Tr} (u (g h ^ {- 1}))   d g.
\]

\begin{enumerate}
\def\labelenumi{\alph{enumi})}
\setcounter{enumi}{1}
\tightlist
\item
  设 \(f\) 是 \(\mathbf{G}\) 上的可积函数；\(f\) 几乎处处等于某个无穷次可微函数，当且仅当函数 \(u \mapsto \| u(f) \|_{\infty}\) 在 \(\hat{\mathbf{G}}\) 上速降。
\end{enumerate}

设 \(f\) 是 \(G\) 上的无穷次可微函数，且 \(\Gamma\) 是 \(G\) 的 Casimir 元；由公式 (14)，

\[
\tilde {\Gamma} (u) ^ {n} u (f) = u (f) u (\Gamma) ^ {n} = u ((\mathrm{L} _ {\Gamma}) ^ {n} f)
\]

对于所有 \(n \geq 0\)，从而

\[
\tilde {\Gamma} (u) ^ {n} \| u (f) \| _ {\infty} \leq \| (\mathrm{L} _ {\Gamma}) ^ {n} f \| _ {1} \leq \sup _ {g \in \mathrm{G}} | ((\mathrm{L} _ {\Gamma}) ^ {n} f) (g) |;\tag{15}
\]

因此，函数 \(u \mapsto \|u(f)\|_{\infty}\) 确实是速降的。

反之，设 \(\mathrm{A} = (\mathrm{A}_{u})_{u \in \hat{\mathrm{G}}}\) 是 \(\mathrm{F}(\hat{\mathrm{G}})\) 的一个元素，且函数 \(u \mapsto \|A_{u}\|_{\infty}\) 是速降的。令 \(f_{u}(g) = \langle u(g)|A_{u} \rangle\)；函数 \(g \mapsto f_{u}(g)\) 是解析的，因而是无穷次可微的。由 Chap. III, §3, no. 7, Prop. 27，

\[
(\mathrm{L} _ {x} f _ {u}) (g) = \langle u (g) u (x) | \mathrm{A} _ {u} \rangle
\]

对于所有 \(x \in \mathfrak{g}\)。设 \(t \in \mathrm{U}(\mathrm{G})\)；由前式，

\[
(\mathrm{L} _ {t} f _ {u}) (g) = \langle u (g) u (t) | \mathrm{A} _ {u} \rangle
\]

从而

\[
\begin{array}{r} | (\mathrm{L} _ {t} f _ {u}) (g) | = | \langle u (g) u (t) | \mathrm{A} _ {u} \rangle \leq d (u) ^ {2} \| u (t) \| _ {\infty} \| u (g) \| _ {\infty} \| \mathrm{A} _ {u} \| _ {\infty} \\ = d (u) ^ {2} \| u (t) \| _ {\infty} \| \mathrm{A} _ {u} \| _ {\infty}. \end{array}
\]

由于 \(d(u)\) 和 \(\|u(t)\|_{\infty}\) 是缓增的（Prop. 3），而 \(\|A_{u}\|_{\infty}\) 是速降的，函数 \(u \mapsto \sup_{g} |(\mathrm{L}_{t}f_{u})(g)|\) 是速降的；因此族 \((\mathrm{L}_{t}f_{u})_{u \in \hat{\mathrm{G}}}\) 一致可求和。由此 \(^{7}\) 可知，族 \((f_{u})\) 的和是 G 上的无穷次可微函数，其 Fourier 余变换为 \((\mathrm{A}_{u})\)，定理得证。

记 \(\mathcal{S}(\hat{\mathrm{G}})\) 为 \(\mathrm{L}^2 (\hat{\mathrm{G}})\) 的向量子空间，它由族 \(\mathrm{A} = (\mathrm{A}_u)_{u\in \hat{\mathrm{G}}}\) 构成，其中函数 \(u\mapsto \| \mathrm{A}_u\|_\infty\) 在 \(\hat{\mathrm{G}}\) 上速降。由定理可知，映射 \(\overline{\mathcal{F}}:f\mapsto (u(f))_{u\in \hat{\mathrm{G}}}\) 和 \(\mathcal{F}\colon \mathrm{A}\mapsto \sum_{u\in \hat{\mathrm{G}}}\langle u(g)|\mathrm{A}_u\rangle\) 诱导出复向量空间 \(\mathcal{C}^\infty (\mathrm{G};\mathbf{C})\) 与 \(\mathcal{S}(\hat{\mathrm{G}})\) 之间的互逆同构。在空间 \(\mathcal{C}^\infty (\mathrm{G};\mathbf{C})\) 上赋予一致 \(\mathrm{C}^\infty\) 收敛的拓扑 (§6, no. 4)，该拓扑可由半范数族 \(f\mapsto \sup_{g\in \mathrm{G}}|\mathrm{L}_t f(g)|\)（\(t\in \mathrm{U}(\mathrm{G})\)）定义；在空间 \(\mathcal{S}(\hat{\mathrm{G}})\) 上赋予由半范数序列 \(p_n:\mathrm{A}\mapsto \sup_{u\in \hat{\mathrm{G}}}(\tilde{\Gamma} (u) + 1)^n\| \mathrm{A}_u\|_\infty\) 定义的拓扑。前一证明中的公式 (15) 表明 \(\overline{\mathcal{F}}\) 是连续的。设 \(t\in \mathrm{U}(\mathrm{G})\)，且 \(\mathrm{A} = (\mathrm{A}_u)_{u\in \hat{\mathrm{G}}}\) 是 \(\mathcal{S}(\hat{\mathrm{G}})\) 的一个元素；令 \(f_{n}(g) = \langle u(g)|\mathrm{A}_{u}\rangle\)。设 \(p\) 是满足 \(\sum_{u\in \hat{\mathrm{G}}} \tilde{\Gamma} (u)^{-p} = \mathrm{M} < \infty\) 的整数。由前一证明，存在正整数 \(m\)，使得对于所有 \(g\in \mathrm{G}\)，

\[
| (\mathrm{L} _ {t} f _ {u}) (g) | \leq d (u) ^ {2} \| u (t) \| _ {\infty} \| \mathrm{A} _ {u} \| _ {\infty} \leq m. (1 + \tilde {\Gamma} (u)) ^ {m} \tilde {\Gamma} (u) ^ {- p} \| \mathrm{A} _ {u} \| _ {\infty}
\]

所以 \(|(\mathrm{L}_t\mathcal{F}(\mathrm{A}))(g)| \leq m\mathrm{M}p_m(\mathrm{A})\)；这证明了 \(\mathcal{F}\) 是连续的。因此：

推论。映射 \(\overline{\mathcal{F}}: f \mapsto (u(f))_{u \in \hat{\mathrm{G}}} \text{ 和 } \mathcal{F}: \mathrm{A} \mapsto \sum_{u \in \hat{\mathrm{G}}} \langle u(g) | \mathrm{A}_u \rangle\) 诱导出拓扑向量空间 \(\mathcal{C}^\infty(\mathrm{G}; \mathbf{C})\) 与 \(\mathcal{S}(\hat{\mathrm{G}})\) 之间的互逆同构。

